% 第六章 第二平台：AMD MI350X
\chapter{第二平台：AMD MI350X}\label{ch:mi350}

第五章的結果都來自同一個 GH200 叢集。為了檢驗結論能否移植到另一家廠商、另一種功耗管理，我們在兩個各有 8 張 AMD Instinct MI350X 的節點上重做主要實驗：機制對照、矩陣乘法、FP64 HPL、HPL-MxP、DF-MP2 與高斯過程，並與 INT8 模擬（GEMMul8）正面比較。本章也用這些量測驗證第 \ref{sec:powermodel} 節的功耗模型。每個實驗都在第二個節點上重複一次；除非特別說明，百分比都相對同一節點、同一功耗上限下未修改的原生程式。

% ==================================================================
\section{平台、量測與統計單位}\label{sec:mi-setup}

兩個節點（n01、n02）各有 8 張 MI350X OAM GPU（gfx950，288~GB HBM3E，最高時脈 2200~MHz），軟體為 ROCm 7.2.3（rocBLAS、hipBLASLt、rocSOLVER、RCCL）與 Open MPI 5.0.7（UCX）。功耗上限是每張 GPU 的封裝功耗上限（PPT），以 \texttt{amd-smi} 設定，可設 1–1000~W；預設 1000~W 也是最高值，不含 CPU。每張 GPU 的閒置功耗為 244–275~W。

與 GH200 不同，MI350X 的 8 張 GPU 同時執行同一個設定，共用機箱的風扇與電源。我們每 100~ms 讀取 8 張 GPU 的功耗、時脈、溫度，以及硬體的能量累加器（每單位 15.26~$\mu$J，能量取視窗兩端的差，不受取樣頻率影響）和功耗、溫度限制的累計時間；每個量測視窗都記錄實際生效的上限。HIP 的裝置編號與 PCI 位址的順序不同，所以分析一律以程式印出的 PCI 位址對應功耗記錄。

同一個機箱、同一個工作中的量測並不獨立，因此本章的信賴區間以真正獨立的單位計算：矩陣乘法以「一個量測視窗」（8 張 GPU 先平均）為單位，HPL 與 HPL-MxP 以「一次整節點執行」為單位，DF-MP2 與高斯過程以「一輪」（8 張 GPU 先平均）為單位，兩個節點的單位合併後取 95\% $t$ 區間。各變體的順序為成對交錯（ABBA，或後述的成對 ABBA）。溫度的影響與控制見第 \ref{sec:mi-temp} 節。

% ==================================================================
\section{預設上限就受限：省電機制與時脈}\label{sec:mi-mech}

MI350X 在出廠的 1000~W 上限下就受功耗限制。原生 FP64 矩陣乘法有 82\% 的時間停在上限，時脈 1890~MHz；保留 0 位時只耗 785~W、時脈 2183~MHz，快 17\%。在我們的量測中，FP64、BF16 與 FP8 矩陣乘法在 1000~W 有 82–100\% 的時間撞上限。GH200 的預設上限下，截斷多半只表現為省電；MI350X 則在出廠設定下就直接轉成速度。

\textbf{位元不翻轉還是數值為零？}與第 \ref{sec:mech} 節相同，我們把被捨去的位元填成全 1 或固定樣式：位元同樣不再翻轉，但數值不是 0。MI350X 可以鎖定最高時脈，所以我們把時脈鎖在 900、1300 或 1700~MHz（低於任何設定會達到的時脈，功耗限制的時間只有 0–0.6\%），直接比較每次運算的能量（8 張 GPU × ABBA）。表~\ref{tab:mi-fill} 與 GH200 並列：在兩家廠商的矩陣引擎、FP64 向量單元與記憶體搬移上，填值都保留了大部分的效益，所以主要原因是位元不再翻轉，而不是數值為 0。這些比例不是精確的因果分解：填值也會改變乘法器內的部分積、對齊與累加。比例幾乎不隨時脈改變（FP64 保留 27 位：900、1300、1700~MHz 分別為 83、81、82\%），總省電量則隨時脈增加（FP64 保留 0 位：每次運算少耗 20、24、29\%），與 $CV^2f$ 相對靜態功耗變大的預期一致。效益也不限於矩陣引擎：FP64 乘加迴圈保留 0 位在 600~W 快 33\%，在 600~W 受限的記憶體複製快 28\%。

\begin{table}[tbp]
\centering
\caption[兩平台的填值對照]{以全 1／固定樣式填入被捨去的位元時，保留的歸零效益比例（\%）。GH200：1000~W 的模組功耗或 600~W 的吞吐量；MI350X：鎖頻下每次運算的能量（矩陣乘法為三個時脈的範圍）、600~W 的吞吐量或 1000~W 的功耗。最後一欄是 MI350X 上歸零本身的效益。}\label{tab:mi-fill}
\small
\begin{tabularx}{\textwidth}{@{}LccL@{}}
\toprule
核心與保留位數 & GH200 & MI350X & MI350X 歸零的效益 \\
\midrule
FP64 矩陣乘法，0 位 & 76/68 & 73–75/69 & 每次運算少耗 20–29\% \\
FP64 矩陣乘法，20 位（GH200）、27 位（MI350X） & 87/82 & 81–83/73–74 & 每次運算少耗 7–10\% \\
BF16 矩陣乘法，0 位 & 92/92 & 77–78/68–72 & 每次運算少耗 13–15\% \\
FP8 矩陣乘法，0 位 & — & 65–66/61 & 每次運算少耗 9–10\% \\
FP64 乘加迴圈，0 位 & 87/62 & 74/63 & 快 33\% \\
記憶體複製，0 位 & 97/86 & 83/41 & 功耗少 9\% \\
\bottomrule
\end{tabularx}
\begin{figsource}鎖頻填值 \path{mi350/logs/e1b_clk900_*}、\path{e1b_clk1300*}、\path{e1b_clk1700_*}；非矩陣乘法核心 \path{mi350/logs/n02/nongemm_*}\end{figsource}
\end{table}

\textbf{鎖頻下仍隨資料而變的時脈。}時脈鎖定、完全不受功耗限制時，有效時脈仍隨資料改變：鎖在 1300~MHz 時，原生 FP64 矩陣乘法穩定在 1171–1179~MHz（8 張 GPU 的第 5 到 95 百分位），保留 0 位則為 1235–1240~MHz；BF16 原生 1088–1100~MHz，保留 0 位 1132–1142~MHz，溫度只有 59–63\,°C。功耗與溫度限制的時間都是 0，每個 XCD 的「時脈低於主機上限」計數器即使在功耗限制生效時也一直是 0，所以沒有任何現有的計數器回報這種降速。電源電壓瞬間下陷時的電路層級時脈延展（clock stretching）可以解釋這個現象，但我們沒有直接量測；截斷同樣減輕了它。第 \ref{sec:powermodel} 節的模型以 $\sigma$ 描述這個效應。

% ==================================================================
\section{功耗上限下的矩陣乘法}\label{sec:mi-gemm}

測試使用 rocBLAS／hipBLASLt 矩陣乘法（FP64 為 $32768\times32768\times2048$，BF16 與 FP8 為 $k=8192$），輸入為均勻隨機值，每個上限以 ABBA 順序量測兩個視窗。歸零在每個上限下都加速（圖~\ref{fig:mi-gemm}）。在 1000~W，FP64 保留 36、27、0 位分別快 5.1、9.2、16.9\%，BF16 保留 3、0 位快 7.7、15.6\%，FP8 保留 0 位快 8.8\%（兩節點共 4 個視窗的平均，95\% 區間 $\pm$0.6–1.6 個百分點）；在 500~W，同樣的設定快 8.6、15.2、60.0、13.5、30.6 與 17.4\%（$\pm$0.1–1.5）。兩個節點相差最多 1.8 個百分點，溫度校正改變任何數值都不超過 1.1 個百分點。FP64 保留 36 位在 600–900~W 相當於每張 GPU 多給 42–72~W。作為對照，GH200 在預設的 1000~W 下 FP64 保留 35 位只快 3.9\%（第 \ref{sec:gemm} 節）。

\begin{figure}[tbp]
\centering
\figfile{fig6_01_gemm_caps}
\caption[MI350X 矩陣乘法在各功耗上限下的加速]{MI350X 矩陣乘法在各封裝功耗上限下保留不同位數的加速。實心：節點 n02 兩個量測視窗的平均（每個視窗先平均同時執行的 8 張 GPU），誤差棒為兩個視窗的範圍；空心：節點 n01。部分系列左右錯開約 8~W 以免重疊。}\label{fig:mi-gemm}
\begin{figsource}\path{mi350/results/e1n02_1029_0502_raw.csv}（n02）、\path{e1_1028_1709_raw.csv}（n01）\end{figsource}
\end{figure}

這些短視窗不是暖機造成的假象：在穩態下（每個設定先執行 4 分鐘再量 2 分鐘），FP64 保留 36 位與 BF16 保留 0 位在 800~W 快 6.1 與 26.4\%、1000~W 快 5.4 與 15.6\%，與同一節點的短視窗（7.6、25.2、5.0、14.8\%）相差不超過 1.5 個百分點（第 \ref{sec:mi-temp} 節）。

% ==================================================================
\section{功耗模型的驗證}\label{sec:mi-model}

第 \ref{sec:powermodel} 節的模型把 $P_0$ 設為實測的閒置功耗（GH200 模組 174~W，MI350X 每張 GPU 262~W），不以驗證資料擬合任何參數。MI350X 以 1300~MHz 鎖頻的量測得到 $R$ 與 $\sigma$（不使用上限掃描的任何一點）；GH200 無法鎖頻，以最低上限的速度推得 $R$（$\sigma=1$，該點不列入評估）。每個平台各評估 42 個「上限 × 格式」的點（FP64 保留 0、12/27、35/36 位，BF16 保留 0、3 位，FP8 保留 0 位）。

\begin{figure}[tbp]
\centering
\figfile{fig6_02_model}
\caption[功耗模型預測與實測的矩陣乘法加速]{功耗模型預測與實測的矩陣乘法加速，$P_0$ 為實測的閒置功耗。每個點是一個功耗上限與格式的組合。}\label{fig:mi-model}
\begin{figsource}\path{mi350/results/model_eval_points.csv}（\path{mi350/analysis/model_eval.py}）\end{figsource}
\end{figure}

\begin{table}[tbp]
\centering
\caption[功耗模型的誤差]{功耗模型預測加速的誤差（百分點），每個平台 42 個點。閒置：$P_0$ 為實測閒置功耗；留一：以其他格式擬合 $P_0$ 後預測被留下的格式；零：把所有功耗都當成動態功耗。}\label{tab:mi-modelerr}
\small
\begin{tabular}{@{}llrrr@{}}
\toprule
平台 & $P_0$ & 平均誤差 & 最大誤差 & FP64 27–36 位的最大誤差 \\
\midrule
GH200 & 閒置 174~W & 1.5 & 4.6 & 2.9 \\
GH200 & 留一（190–200~W） & 1.0 & 3.6 & — \\
GH200 & 0 & 6.1 & 12.9 & — \\
MI350X & 閒置 262~W & 2.1 & 10.4 & 2.9 \\
MI350X & 留一（190–250~W） & 2.3 & 9.4 & — \\
MI350X & 0 & 4.0 & 17.2 & — \\
\bottomrule
\end{tabular}
\end{table}

平均誤差在 GH200 為 1.5 個百分點（最大 4.6），在 MI350X 為 2.1（最大 10.4；圖~\ref{fig:mi-model}、表~\ref{tab:mi-modelerr}）。若改以其他格式擬合 $P_0$ 再預測被留下的格式，誤差為 1.0 與 2.3；把全部功耗都當成動態功耗（$P_0=0$）則增為 6.1 與 4.0；MI350X 若省略 $\sigma$ 為 2.5。實測的比值為：FP64 保留 36、27、0 位的 $R=0.94$、0.89、0.67，BF16 保留 3、0 位為 0.92、0.85，FP8 保留 0 位為 0.90，$\sigma=1.01$–1.05。應用實際使用的 FP64 寬度（27–36 位）在兩個平台上誤差都不超過 2.9 個百分點；誤差最大的是 MI350X 上 BF16 保留 0 位在 400~W（原生幾乎跑不動），以及 FP64 保留 0 位在 500–800~W（時脈超出原生曲線的範圍）。

\textbf{外部驗證。}上述矩陣乘法的點參與了模型的發展。為了在沒有參與的程式上檢驗，我們在節點 n02 上以 1300~MHz 鎖頻執行 DF-MP2 與高斯過程，從整個程式的功耗與時脈得到 $R$ 與 $\sigma$（DF-MP2 保留 32、20 位：$R=0.92$、0.86，$\sigma=1.01$、1.02；高斯過程保留 36、28 位：$R=0.96$、0.93），再由同一節點上原生程式的效能曲線預測 500–1000~W 的加速。16 個點的平均誤差為 1.6 個百分點，最大 3.6。DF-MP2 的誤差在 $-1.6$ 到 $+2.4$ 之間；高斯過程在每個上限都被低估 0.6–3.6 個百分點，可能是因為它的 Cholesky 交替著以矩陣乘法為主與非矩陣乘法的階段，兩者功耗不同，單一個整體比值無法描述。因此只要知道平台的閒置功耗與應用的原生曲線，一次鎖頻量測就能把截斷版本的效能曲線預測到幾個百分點以內；$\Delta W\propto P-P_0$ 也說明了為什麼等效功耗隨上限增加，而相對加速並不會。

% ==================================================================
\section{HPL 與 HPL-MxP}\label{sec:mi-hpl}

\textbf{安全的位數因平台而異。}表~\ref{tab:mi-hplres} 列出 MI350X 單節點 8 張 GPU 在三種問題大小下的殘差。28 位在每個大小都失敗（超過門檻 6–32 倍），30 位也是；32 位勉強通過；36 位有 35–100 倍的餘裕。GH200 上 28 位在四個大小中有三個通過（第 \ref{sec:hpl} 節）。兩個平台的 HPL 版本、DGEMM 的分塊與累加順序都不同，所以在一個平台驗證過的位數不能直接搬到另一個平台。

\begin{table}[tbp]
\centering
\caption[MI350X 上 HPL 的殘差]{MI350X（一個節點 8 張 GPU，1000~W）上 HPL 的殘差（門檻 16，重複執行完全相同）。位數為含隱藏位的有效位數，最近捨入（平手取偶）；(F) 表示未通過檢查。}\label{tab:mi-hplres}
\small
\begin{tabular}{@{}rrrrrr@{}}
\toprule
$N$ & 原生 & 36 位 & 32 位 & 30 位 & 28 位 \\
\midrule
262144 & \ee{5.4}{-5} & 0.46 & 11.1 & 84 (F) & 514 (F) \\
393216 & \ee{4.6}{-5} & 0.15 & 12.0 & 20.2 (F) & 104 (F) \\
540672 & \ee{7.8}{-6} & 0.27 & 3.5 & 17.3 (F) & 102 (F) \\
\bottomrule
\end{tabular}
\begin{figsource}\path{mi350/logs/hplbits262k_*}、\path{hplbits393k_*}、\path{hplbits540k_*}\end{figsource}
\end{table}

\textbf{FP64 HPL。}原生 HPL（$N=540672$，$1\times8$）在 1000~W 兩個節點分別為 424.9 與 429.8~TFLOP/s。保留 36 位在預設上限快 4.6–5.1\%，800~W 快 4.5–4.9\%，650~W 快 6.2–8.0\%，500~W 快 6.3–7.3\%（兩節點合併：1000~W $4.8\pm0.3$\%、500~W $7.0\pm0.6$\%，各 6 次執行；表~\ref{tab:mi-robust}），每次求解少耗能 4.6–6.8\%，在 650–800~W 的等效功耗為每張 GPU 39–50~W（每個節點約 0.35–0.4~kW）。

\textbf{HPL-MxP。}原生 HPL-MxP（$N=1034752$，$4\times2$，對角優勢的測試矩陣）在 1000~W 為 7.61 與 7.65~PFLOP/s，保留 3 位時為 7.98 與 8.04~PFLOP/s（+4.9 與 +5.1\%），在 500–800~W 快 9.2–19.5\%，每次求解少耗能 5.3–14.6\%。迭代修正從 2 次增為 4 次，約佔 LU 時間的 1.5\%。保留 0 位在 MI350X 上不收斂，所以兩個平台的操作點都是 3 位。GH200 上的 HPL-MxP 在 1000~W 會觸發機箱的功耗煞車而不報；MI350X 沒有這個問題。這些結果只適用於對角優勢的測試矩陣：程式的產生器一律把對角元設為該列絕對值的和，我們沒有正式規則允許的產生器，所以 3 位的操作點與其收益都不是 HPL-MxP 的一般性質。不過它們不依賴對角 tile 的近似分解：改用精確分解（\texttt{-{}-exact-diag}，節點 n02，一組 ABBA），保留 3 位仍通過殘差檢查（0.43），在 1000~W 快 4.4\%、650~W 快 9.4\%（溫度校正後 4.4 與 9.3\%），近似分解則為 5.1 與 9.7\%。

\begin{figure}[tbp]
\centering
\figfile{fig6_03_apps}
\caption[MI350X 上四個應用在各上限下的加速]{MI350X 上程式內截斷相對未修改程式的加速，兩節點合併。誤差棒為獨立單位的 95\% $t$ 區間：HPL 與 HPL-MxP 為整節點的一次執行，高斯過程與 DF-MP2 為 8 張 GPU 的一輪；$n$ 列於圖例。各節點與溫度校正後的數值見表~\ref{tab:mi-robust}。}\label{fig:mi-apps}
\begin{figsource}\path{mi350/results/apps_units.csv}（\path{mi350/analysis/apps_units.py}）\end{figsource}
\end{figure}

\begin{table}[tbp]
\centering
\caption[MI350X 的主要結果（兩節點、獨立單位）]{MI350X 在預設上限與 500~W 的加速（\%）。合併：兩節點所有獨立單位的平均 $\pm$ 95\% $t$ 區間（視窗：8 張 GPU 平均的一個量測視窗；執行：整節點的一次執行；輪：8 張 GPU 的一輪）。n01、n02：各節點的原始值（相同散熱環境）／校正到相同 junction 溫度的值。}\label{tab:mi-robust}
\footnotesize\setlength{\tabcolsep}{3pt}
\begin{tabular}{@{}llcccccc@{}}
\toprule
& & \multicolumn{3}{c}{1000~W} & \multicolumn{3}{c}{500~W} \\
\cmidrule(lr){3-5}\cmidrule(l){6-8}
結果 & 單位，$n$ & 合併 & n01 & n02 & 合併 & n01 & n02 \\
\midrule
矩陣乘法 FP64，36 位 & 視窗，4 & $5.1\pm0.6$ & 5.2/5.4 & 5.0/5.5 & $8.6\pm0.1$ & 8.6/8.6 & 8.6/8.7 \\
矩陣乘法 BF16，0 位 & 視窗，4 & $15.6\pm1.4$ & 16.4/16.6 & 14.8/15.2 & $30.6\pm1.5$ & 29.8/29.4 & 31.4/31.6 \\
矩陣乘法 FP8，0 位 & 視窗，4 & $8.7\pm1.6$ & 9.6/9.5 & 7.9/7.9 & $17.4\pm0.6$ & 17.1/17.2 & 17.7/17.7 \\
HPL，36 位 & 執行，6 & $4.8\pm0.3$ & 4.6/4.5 & 5.1/4.8 & $7.0\pm0.6$ & 7.3/7.2 & 6.3/6.5 \\
HPL-MxP，3 位 & 執行，8 & $5.0\pm0.4$ & 4.9/4.8 & 5.1/5.1 & $19.3\pm1.3$ & 19.5/18.8 & 19.1/21.3 \\
高斯過程，28 位 & 輪，4 & $7.4\pm0.6$ & 7.6/7.1 & 7.1/7.3 & $11.0\pm1.0$ & 10.8/12.5 & 11.2/12.9 \\
DF-MP2，20 位 & 輪，4 & $10.3\pm0.6$ & 10.2/13.3 & 10.4/14.2 & $18.6\pm0.3$ & 18.8/19.8 & 18.5/20.0 \\
DF-MP2，32 位 & 輪，4 & $5.7\pm0.4$ & 5.5/7.5 & 5.9/7.4 & $9.8\pm0.1$ & 9.7/10.5 & 9.8/11.0 \\
\bottomrule
\end{tabular}
\begin{figsource}\path{mi350/results/apps_units.csv}、\path{*_tempadj.csv}\end{figsource}
\end{table}

% ==================================================================
\section{DF-MP2 與高斯過程}\label{sec:mi-apps}

\textbf{DF-MP2。}GPU4PySCF 不能在 MI350X 上執行，所以我們以 PySCF 在 CPU 上建立相同的三指標張量（taxol，def2-TZVP，RI 輔助基底），再以自寫的 HIP 程式在 GPU 上收縮：每一對佔據軌域 $(i, j\le i)$ 以一次 FP64 DGEMM 計算 $B_iB_j^{\mathsf T}$，內積長度為輔助基底的大小（$k=5477$）。相關能與 PySCF 相差 \ee{2}{-13}\Eh 以內，與 GPU4PySCF 相差 \ee{7.5}{-8}\Eh。張量建好後原地歸零一次（最近捨入）：保留 32 位時 MP2 相關能最多改變 \ee{5}{-12}\Eh，20 位最多 \ee{2.2}{-8}\Eh。在 1000~W 原生有 95\% 的時間撞上限，原生時間在 1000、800、650、500~W 分別為 21.7、26.4、34.1、55.8~s。兩節點共 4 輪，保留 32 位快 $5.7\pm0.4$、$6.7\pm1.4$、$9.5\pm0.7$、$9.8\pm0.1$\%，保留 20 位快 $10.3\pm0.6$、$12.7\pm1.8$、$16.6\pm0.9$、$18.6\pm0.3$\%。GH200 在 1000~W 只快 0.5–0.8\%（主要表現為省電，第 \ref{sec:mp2} 節）；MI350X 在預設上限就快 5.7–10.3\%。

\textbf{高斯過程。}我們把第 \ref{sec:gp} 節的似然函數移植為 HIP 程式（rocBLAS／rocSOLVER 的分塊 Cholesky，$n=80000$，Matérn-3/2，每個面板一次 FP64 DGEMM 做尾端更新，$k=1024$），在每個面板做尾端更新前原地截斷（最近捨入）。對數似然約為 $1.2$–$1.7\times10^{5}$，保留 36 位最多偏離 \ee{1.9}{-5}，28 位最多 \ee{5.3}{-3}，相對誤差約 \ee{3}{-8}，遠小於估計所需。兩節點共 4 輪，保留 36 位在 1000、800、650、500~W 快 $3.9\pm0.2$、$4.1\pm0.6$、$5.3\pm1.7$、$6.9\pm0.7$\%，保留 28 位快 $7.4\pm0.6$、$7.1\pm1.8$、$7.8\pm2.8$、$11.0\pm1.0$\%。

% ==================================================================
\section{與 INT8 模擬（GEMMul8）的正面比較}\label{sec:mi-emu}

MI350X 的 INT8 與 FP64 吞吐量比約為 GH200 的兩倍，對 INT8 模擬較有利。表~\ref{tab:mi-h2h} 讓兩種方法各用自己允許的最佳設定並列比較：在內積長、上限寬鬆、模擬的固定成本能攤銷時（HPL、DF-MP2、$k\ge2048$ 的矩陣乘法），模擬大幅領先；上限收緊時領先縮小，甚至反轉；內積短時（高斯過程、$k=1024$ 的矩陣乘法）模擬比原生還慢。截斷的收益較小，但處處為正、隨上限收緊而增加、不需要額外記憶體，而且適用於沒有模擬可用的 BF16 與 FP8。

\begin{table}[tbp]
\centering
\caption[MI350X 上截斷與 GEMMul8 的正面比較]{MI350X 上截斷與 GEMMul8 相對最佳原生的加速：預設上限／500~W，括號內為誤差。截斷：HPL 36 位、矩陣乘法 35 位、DF-MP2 20 位、高斯過程 28 位；GEMMul8：HPL 最佳設定 10 個模數，其餘 12 個。第二列 HPL 在同一設定、相近殘差下比較。矩陣乘法 $m=n=32768$。}\label{tab:mi-h2h}
\small
\begin{tabularx}{\textwidth}{@{}Lll@{}}
\toprule
工作負載 & 截斷 & GEMMul8 \\
\midrule
FP64 矩陣乘法，$k=1024$ & +5.1/+9.2\%（\ee{8}{-12}） & $-13$/$-45$\%（\ee{4}{-13}） \\
FP64 矩陣乘法，$k=4096$ & +5.5/+9.0\%（\ee{8}{-12}） & +84/+18\%（\ee{8}{-13}） \\
HPL，各自的最佳設定（殘差） & +4.6/+7.3\%（0.27） & +36/+11\%（1.79） \\
HPL，$N=458752$，$4\times2$（殘差） & +3.7/+8.0\%（0.086） & +32/$-0.4$\%（0.033） \\
DF-MP2（$|\Delta E|$，\Eh） & +10/+19\%（\ee{2}{-8}） & +67/+34\%（\ee{7}{-12}） \\
高斯過程，$n=64000$（$|\Delta\ell|$） & +5.9/+11.8\%（\ee{2}{-3}） & $-11$/$-43$\%（\ee{3}{-7}） \\
HPL-MxP（BF16） & +4.9/+19\%（殘差 0.5） & 沒有對應的模擬 \\
\bottomrule
\end{tabularx}
\end{table}

\textbf{分界：形狀與上限。}圖~\ref{fig:mi-xover} 是 $m=n=32768$ 時兩種方法相對原生 FP64 的加速隨內積長度 $k$ 的變化。截斷到 35 位（誤差 \ee{7.8}{-12}）在每個 $k$ 都比原生快，1000~W 快 5–6\%、500~W 快 8–9\%。GEMMul8 用 12 個模數（較精確，約 \ee{5}{-13}）在 1000~W 時，$k=1024$ 比原生慢 13\%，$k=2048$、4096、8192 則快 36、84、121\%；在 500~W 變為 $-45$、$-18$、+18、+48\%。$k=2048$ 時，模擬在 1000 與 700~W 勝出，截斷在 500~W 勝出。模擬一直用滿功耗上限，所以受限時額外的整數運算要以時脈支付；截斷則是從功耗預算中拿掉工作。GH200 上的 GEMMul8 也是如此，而只用約 706~W 的 cuBLAS 模擬在低上限反而相對有利（第 \ref{sec:crossover} 節）：方向取決於模擬本身的功耗。

\begin{figure}[tbp]
\centering
\figfile{fig6_04_crossover}
\caption[MI350X 上截斷與 GEMMul8 的分界]{MI350X 上截斷（最近捨入）與 GEMMul8 相對原生 FP64 的加速隨內積長度 $k$ 的變化（$m=n=32768$，節點 n02，8 張 GPU × ABBA），三個封裝功耗上限。上限收緊時，模擬的曲線全部下降，截斷的曲線全部上升。}\label{fig:mi-xover}
\begin{figsource}\path{mi350/results/xover_1029_1508_summary.csv}\end{figsource}
\end{figure}

\textbf{應用。}DF-MP2（$k=5477$）中，GEMMul8 用 12 個模數比原生快 34–67\%，每個上限都勝過截斷，但對保留 20 位的領先從 1000~W 的 57 個百分點縮小到 500~W 的 16 個百分點（GH200 上的 cuBLAS 模擬反而比原生慢 18–35\%）。FP64 HPL 中，GEMMul8 用 10 個模數、在它自己的最佳設定（$N=491520$，NB 4096，$4\times2$；殘差 1.79，36 位截斷為 0.03）在 1000~W 達到 576.8~TFLOP/s，比最佳原生設定高 36\%，但對最佳截斷的領先從 1000~W 的 30\% 縮小到 500~W 的 4\%。在相近的精度下，上限會反轉勝負：用 12 個模數（殘差 0.033，36 位截斷為 0.086；$N=458752$、$4\times2$，是放得下的最大問題），GEMMul8 在 1000、800、650~W 比原生快 32.0、26.5、14.9\%，在 500~W 卻是 $-0.4$\%，此時截斷快 8.0\%（溫度校正後為 $-1.5$ 與 8.0\%）。高斯過程的尾端更新 $k=1024$，GEMMul8 在每個上限都比原生慢（$n=64000$ 時結果正確：12 個模數在 1000、800、650、500~W 為 $-11$、$-22$、$-37$、$-43$\%，對數似然差 \ee{3}{-7} 以內），保留 28 位的截斷則快 5.9–11.8\%。

\textbf{只在應用中出現的成本。}模擬需要工作空間：最佳原生 HPL 的 $N=540672$ 時，矩陣本身就佔每張 GPU 292~GB，GEMMul8 的工作空間放不下；$1\times8$ 的格網即使在 $N=458752$ 也要求 100–114~GB；12 個模數連 $N=491520$ 都失敗，HPL 必須先縮小問題才能使用模擬。每遇到一種新的矩陣形狀，GEMMul8 都有一次很大的初始化成本：DF-MP2 第一次走過 226 種形狀要 142~s，而不是 13~s（上述時間已扣除這次暖機）。

\textbf{靜默的錯誤。}高斯過程在 $n=80000$ 時，第一次尾端更新（$m=n=78976$，$k=1024$）被攔截的 DGEMM 回報成功，結果卻是錯的，256 次執行無一例外，也沒有任何訊息。我們以均勻隨機輸入寫了一個最小重現程式（\path{mi350/apps/g8_repro.hip}），精確定出界線：不論直接呼叫函式庫還是透過攔截，$k=1024$ 的 $C=AB^{\mathsf T}$ 在 $m=n=65535$ 時正確（相對誤差 \ee{1.6}{-15}），從 $m=n=65536$ 起就錯（相對誤差 31，所有元素都是有限值），也就是 $mn$ 達到 $2^{32}$ 時，與我們所用的 GEMMul8 3.2.0 HIP 版本中 32 位元大小溢位的情況一致。這是實作上的缺陷，與資料無關，也不是模擬方法本身的性質；本報告其他的模擬結果都在這個大小以下。它說明了整個替換 BLAS 呼叫的風險：在一個大小驗證過的設定，換一個大小就可能安靜地出錯。

\textbf{指數跨度。}與 GH200 相同（第 \ref{sec:span} 節），均勻合成的指數散佈幾乎不影響 GEMMul8：每個元素或每列的指數散佈到 64 個 binade，12 個模數的逐元素誤差中位數仍為 \ee{5.7}{-13}–\ee{8.0}{-13}，第 99 百分位最多 \ee{4.4}{-11}。GH200 上的失效證據來自真實資料（纈胺黴素），列內跨度寬只是必要條件。

\textbf{在模擬之內再截斷。}先截斷輸入再交給 GEMMul8，希望降低 INT8 引擎的翻轉：在受限時只多 1–4 個百分點（12 個模數、19 位輸入，500~W 從快 18.4\% 變為 22.6\%），在 1000~W 沒有差別，誤差卻變成截斷的誤差（\ee{5.1}{-7}）；不截斷、改用 8 個模數更快也更準。在模擬之內，該調的是模數的個數，不是輸入的位數。

% ==================================================================
\section{量測陷阱：風扇與溫度}\label{sec:mi-temp}

在 MI350X 的 800~W 下，一次三個變體（原生、4 位、3 位）的 ABBA 序列讓 HPL-MxP 保留 4 位看起來比 3 位快（+11.3 對 +9.3\%）。原因是機箱風扇：一次高溫的執行之後，風扇加速，junction 溫度在 20~s 內下降 25\,°C，降溫條件立刻滿足，下一次執行便低溫約 12\,°C；排在中間的變體永遠跟在高溫執行之後。改成每次只比兩個變體的成對 ABBA，並規定最短休息時間，只減輕而沒有消除這個效應：高溫執行之後的那次仍可能低溫 15\,°C、快 4.5\%，因為較冷的晶片漏電較少，留給切換的功耗預算較多。

溫度因此以兩種方式出現，不能混為一談。前一次執行與風扇造成的差異是干擾因素；截斷的程式本身耗電較少而較冷，則是它的收益之一（只在不受功耗限制時；受限時兩者都吃滿上限）。我們報告兩個量：\textbf{相同散熱環境下的收益}是原始的成對比值，由機箱自己控制風扇，是使用者實際得到的；\textbf{相同溫度下的資料效應}則在每個上限內，以變體與該次執行的平均 junction 溫度對效能的對數做迴歸（係數為每 \,°C $-0.1$ 到 $-0.9$\%），移除了干擾，但在截斷本身會降溫時也移除了部分收益。以溫度校正後，被偏誤的那一格中 3 位為 +9.0\%（成對重跑 +9.1\%，n02 為 +9.1\%），4 位為 +4.4–6.2\%；其他 HPL 與 HPL-MxP 的格子校正後最多移動 0.7 個百分點。高斯過程與 DF-MP2 校正後的收益比原始值大 0–4 個百分點：它們的變體在同一個程序中依序執行（原生、32/36 位、20/28 位，再反向），截斷的變體位在每一輪中間最熱的部分，所以我們報告的原始收益是保守的。

相同的平均溫度不代表相同的風扇轉速或熱歷史，所以我們在兩個關鍵設定上直接控制熱歷史，兩者都重現了校正後的數值：
\begin{itemize}
  \item \textbf{穩態}（節點 n02）：每個設定先執行 4 分鐘矩陣乘法、再量 2 分鐘，各自達到穩態。FP64 保留 36 位與 BF16 保留 0 位在 800~W 快 6.1 與 26.4\%，1000~W 快 5.4 與 15.6\%，與短視窗（7.6、25.2、5.0、14.8\%）相差 1.5 個百分點以內。
  \item \textbf{相同的熱歷史}（節點 n01，800~W，成對 ABBA，每個變體 4 次）：每次 HPL-MxP 執行前，8 張 GPU 都先執行同樣 3 分鐘的原生 BF16 矩陣乘法。3 位快 $10.4\pm2.8$\%、4 位快 $6.0\pm3.5$\%（原始），校正後為 9.0 與 6.2\%，與先前各工作的校正值（9.0–9.1 與 4.4–6.2\%）相同；原始數據的排序也恢復正常。
\end{itemize}
我們沒有權限固定風扇轉速，這仍是限制之一（第 \ref{sec:limits} 節）。

% ==================================================================
\section{小結：兩個平台的異同}\label{sec:mi-summary}

表~\ref{tab:mi-cmp} 並列兩個平台的主要結果。機制（翻轉）、應用的排序與「需要以多組大小驗證位數」在兩個平台上一致；不同的是預設上限是否受限、安全的位數，以及與 INT8 模擬的分界。MI350X 的預設上限就壓住時脈，所以截斷在出廠設定下就快 4.6–10\%；FP64 相對 INT8 較弱，使模擬在內積長、上限寬鬆時大幅領先，但在嚴格上限、內積短、記憶體吃緊時仍是截斷較好。

\begin{table}[tbp]
\centering
\caption[GH200 與 MI350X 的比較]{GH200 與 MI350X 的主要結果比較。GH200 的 HPL、HPL-MxP 為 8 個節點；MI350X 為一個節點的 8 張 GPU。}\label{tab:mi-cmp}
\small
\begin{tabularx}{\textwidth}{@{}>{\raggedright\arraybackslash}p{3.3cm}LL@{}}
\toprule
項目 & GH200 & MI350X \\
\midrule
預設上限是否受限 & 部分時間（FP64 矩陣乘法 95\%） & 是（矩陣乘法 82–100\%） \\
填值保留的效益 & 62–97\% & 41–83\% \\
FP64 矩陣乘法 35/36 位，1000~W & +3.9\% & $+5.1\pm0.6$\% \\
HPL 36 位 & 500–800~W：+3.3–4.7\% & 1000~W：$+4.8\pm0.3$\%；500~W：$+7.0\pm0.6$\% \\
HPL 安全位數 & 28 位在 4 個大小中通過 3 個 & 28、30 位全部失敗；36 位安全 \\
HPL-MxP 3 位 & 500~W：+5.8\%；600–800~W：+1.7–2.3\% & 1000~W：$+5.0\pm0.4$\%；500~W：$+19.3\pm1.3$\% \\
DF-MP2 20 位 & 500–700~W：+7.4–7.7\%；1000~W：+0.8\% & 1000~W：$+10.3\pm0.6$\%；500~W：$+18.6\pm0.3$\% \\
高斯過程 28 位 & 500–700~W：+7.5–7.9\%；1000~W：+3.8\% & 1000~W：$+7.4\pm0.6$\%；500~W：$+11.0\pm1.0$\% \\
功耗模型平均誤差 & 1.5 個百分點 & 2.1（兩個應用：1.6） \\
INT8 模擬 & cuBLAS 模擬多半較慢；GEMMul8 只在長內積勝出 & GEMMul8 在長內積、寬鬆上限大幅領先；嚴格上限、短內積時截斷較好 \\
\bottomrule
\end{tabularx}
\end{table}
